% !TeX program = xelatex
\documentclass[12pt,a4paper]{article}
\usepackage{fontspec}
\setmainfont{Times New Roman}
\usepackage[vietnamese]{babel}
\usepackage{amsmath,amssymb}
\usepackage[margin=2.3cm]{geometry}
\usepackage{enumitem}
\setlist[enumerate]{leftmargin=*,itemsep=0.4em}
\setlength{\parindent}{0pt}
\begin{document}
\begin{center}
{\Large\bfseries ĐỀ TOÁN TUYỂN SINH VÀO LỚP 10 THPT HÀ NỘI -- NĂM HỌC 2025--2026}\\[2pt]
{\large\bfseries PHẦN ĐẠI SỐ VÀ TOÁN ỨNG DỤNG}\\
Môn Toán -- Ngày thi: 08/06/2025
\end{center}

\textit{Các câu hình học chứng minh được lược bỏ.}

\section*{Câu I. Thống kê và xác suất}
\begin{enumerate}
\item Kết quả khảo sát 300 học sinh lớp 9 về thời gian tự học của mỗi bạn trong một tuần (đơn vị: giờ) được cho trong bảng:
\[
\begin{array}{c|ccccc}
\text{Thời gian tự học} &[0;4)&[4;8)&[8;12)&[12;16)&[16;20)\\\hline
\text{Số học sinh}&17&72&94&75&42
\end{array}
\]
Xác định tần số và tần số tương đối của nhóm $[12;16)$.
\item Một hộp có 8 chiếc thẻ cùng loại, mỗi thẻ ghi một số trong các số $1,2,3,4,5,6,7,8$. Rút ngẫu nhiên một thẻ. Tính xác suất của biến cố $A$: ``Số ghi trên thẻ rút được chia hết cho 3''.
\end{enumerate}

\section*{Câu II. Biểu thức đại số}
Cho hai biểu thức
\[
A=\frac{\sqrt{x}+2}{\sqrt{x}-2},\qquad
B=\frac{x+\sqrt{x}-4}{x-2\sqrt{x}}-\frac{1}{\sqrt{x}-2},
\quad x>0,\ x\ne4.
\]
\begin{enumerate}
\item Tính giá trị của $A$ khi $x=9$.
\item Chứng minh rằng $B=\dfrac{\sqrt{x}+2}{\sqrt{x}}$.
\item Tìm số nguyên dương $x$ lớn nhất để $\dfrac AB<\dfrac12$.
\end{enumerate}

\section*{Câu III. Bài toán thực tế và hệ thức Viète}
\begin{enumerate}
\item Một ô tô đi từ Hà Nội đến Hải Phòng với vận tốc trung bình $60$ km/h. Khi trở về trên cùng quãng đường, ô tô đi với vận tốc trung bình $40$ km/h. Thời gian đi từ Hà Nội đến Hải Phòng ít hơn thời gian về là 1 giờ. Tính độ dài quãng đường.
\item Bạn Quốc mua một chiếc ba lô và một chiếc máy tính cầm tay có tổng giá niêm yết là 885 nghìn đồng. Cửa hàng giảm giá ba lô 20\% và máy tính cầm tay 25\%; Quốc phải trả 682 nghìn đồng. Tính giá niêm yết của mỗi sản phẩm.
\item Phương trình $x^2+8x-6=0$ có hai nghiệm $x_1,x_2$. Tìm tất cả giá trị của $m$ thỏa mãn
\[
\frac{70-mx_1^2}{x_2}=x_1+mx_2.
\]
\end{enumerate}

\section*{Câu IV. Diện tích và thể tích}
Gia đình bạn Khánh sử dụng một thùng đựng nước dạng hình trụ có bán kính đáy 50 cm, chiều cao 150 cm. Thùng đặt thẳng đứng. Lấy $\pi\approx3{,}14$ và bỏ qua chiều dày của thùng.
\begin{enumerate}
\item Tính diện tích xung quanh của thùng.
\item Sau một thời gian sử dụng, mực nước trong thùng hạ thấp 40 cm. Tính thể tích nước gia đình bạn Khánh đã sử dụng.
\end{enumerate}

\section*{Câu V. Bài toán tối ưu}
Một công ty vận tải đang vận hành 35 xe chở hàng cùng loại. Lợi nhuận trung bình của mỗi xe là 1 triệu đồng một ngày. Cứ bổ sung một xe cùng loại thì lợi nhuận trung bình mỗi xe trong cả đội giảm 20 nghìn đồng một ngày. Hỏi công ty nên bổ sung bao nhiêu xe để tổng lợi nhuận trung bình mỗi ngày của đội xe lớn nhất?
\end{document}
